\subsection{时变碳强度与充电时刻分析}

\subsubsection{不同充电安排对比分析}

电网碳强度采用图\ref{fig:experiment-carbon-intensity}所示的TVGCI-CMDMF-DVRP算例逐小时数据\cite{ref:82}，
日内取值介于0.154至0.644 kgCO$_2$e/kWh之间，峰谷比达4.18倍。
已有研究表明，按时变碳强度重排充电时刻可显著降低电动车充电环节的碳排放\cite{ref:52}；
将其纳入路线与充电的联合优化目标后，路线与充电安排随之改变，排放较不考虑碳强度的方案下降16.6\%--21.4\%\cite{ref:miyabe}。
为考察时变碳强度对降本降碳的作用，在基准条件（北京现行分时电价、基准单位碳价0.20元/kgCO$_2$e）下设置三种充电安排，
对应方案的总成本、各项成本、行驶距离、充电电量与碳排量如表\ref{tab:carbon-charging}所示。
其中，有可用时段即充电指车辆一回场即开始补电、不考虑电价与碳强度，首趟出车前的补电自前一日回场时刻起尽早进行，车型、路线与充电时刻按该策略优化；
考虑时变碳强度（路线不变）指取前者求得的配送路线与车型不变，仅按电价与电网碳强度重排充电时刻；
考虑时变碳强度（路线重优化）指车型、路线与充电时刻按电价与电网碳强度一并重新优化。
下文分别简称为即充方案、路线不变方案与路线重优化方案。
即充方案与路线重优化方案的结果均基于10次运算；路线不变方案由即充方案的10组路线逐组重排得到，
无可行充电窗口的路线沿用即充方案的充电时刻。
碳排量分燃油车直接排放与电动车充电排放两项列出，以便区分车型改派与充电时刻各自的减排。

%%TABLE10%%

各次运算均完成50个客户和13264 kg需求。由表\ref{tab:carbon-charging}可知：
1）路线与车型不变时，考虑时变碳强度降低了成本，却增加了排放。
充电成本平均下降36.75元，总成本平均下降34.50元（1.29\%），10组路线全部下降；
而电动车充电排放平均增加11.28 kgCO$_2$e，10组路线中9组上升。
充电电量在两种安排下逐位相同，这是因为补电量由路线决定、与充电时刻无关，
充电成本的下降全部来自时段单价的差异，排放的上升则来自补电时刻向碳强度更高的时段移动。
2）路线与车型随之调整的额外收益在基准条件下没有显现。
路线重优化方案与路线不变方案的总成本相差3.42元（0.13\%），几乎相同，差值远小于路线搜索10次结果的极差（50.34元）；
路线重优化方案10次中有8次派出3辆电动车，即充方案为5次，其燃油车直接排放比即充方案低22.63 kgCO$_2$e，
但电动车充电排放相应增加，总排放与即充方案几乎相同。
3）在基准条件下，时变碳信号被分时电价信号压制：碳成本占总成本仅1.46\%，考虑时变碳强度的收益全部来自电费，
碳排量不降反升，车型指派与路线也未改变。究其原因，在于分时电价、电网碳强度与作业班次三者在时段上相互错位，
下一小节予以分析。

\subsubsection{分时电价与电网碳强度时段分析}

为分析上述错位，图\ref{fig:carbon-charging}以表\ref{tab:carbon-charging}即充方案的一组路线为例，
给出该路线在两种充电安排下逐时段充电量与电网碳强度、分时电价的对应关系，灰带为企业的两段作业班次。

%%FIGURE4%%

由图\ref{fig:carbon-charging}可知：
1）分时电价与电网碳强度反向。北京谷段电价00:00--07:00的电网碳强度为0.585--0.644 kgCO$_2$e/kWh，是全天最高；
碳强度最低的13:00--15:00（0.154 kgCO$_2$e/kWh）落在平段电价，车辆傍晚回场的15:00--18:00碳强度为0.18--0.47 kgCO$_2$e/kWh，落在平段与峰段电价。
2）低碳价下碳信号无法阻止补电流向谷段。即充方案的三次首趟前补电（87.97 kWh，占全日充电量的39\%）在前一日15:34、16:54与18:10回场时进行，
碳强度分别为0.18、0.27与0.47 kgCO$_2$e/kWh；考虑时变碳强度后全部推至当日06:00的谷段，电价由0.836--1.149元/kWh降至0.563元/kWh，
碳强度却升至0.585 kgCO$_2$e/kWh。以46.94 kWh的一次补电计，电费节省12.82元而碳排量增加18.89 kgCO$_2$e，
按0.20元/kgCO$_2$e折算仅3.78元，考虑碳强度的充电安排仍选择谷段。
单位碳价升至1.00元/kgCO$_2$e后，碳强度为0.18与0.27 kgCO$_2$e/kWh的两次补电不再移动，只有碳强度0.47 kgCO$_2$e/kWh的补电仍移入谷段。
3）作业班次决定了各次补电能够到达的时段。首趟前补电只能在前一日回场到次日开工之间进行，
午休回场的三次补电由09:48--10:57推至12:19--13:00，碳强度由0.24--0.32 kgCO$_2$e/kWh降至0.15--0.16 kgCO$_2$e/kWh，
是考虑时变碳强度真正减排的部分；下午班趟间的两次补电时刻不变。
4）把谷段挪到午间不能改变首趟前补电的去向。车辆回场已在15:00之后，午间谷段不在首趟前补电的可行窗口内，
该部分补电仍移入夜间谷段；因此在单位碳价0.20元/kgCO$_2$e下，谷段挪至午间后路线不变方案相对即充方案的总成本下降71.73元，
碳排量仍增加10.21 kgCO$_2$e；只有单位碳价升至1.00元/kgCO$_2$e后，两者才同时下降（73.64元、4.38 kgCO$_2$e，10组路线全部同时下降）。
5）公共充电站在基准条件下不能成为低碳出口。算例中97个公共充电站与车场功率相同、碳强度相同，客户到最近公共站的路网距离中位数为2.89 km，
但公共站充电须另付0.4元/kWh的服务费，两种充电安排的全部充电会话中在公共站充电的不足1\%。
被锁在夜间的首趟前补电若改在午后到客户附近的公共站补电，每kWh多付0.67元电费、少排0.45 kgCO$_2$e，
盈亏平衡的单位碳价约为1.48元/kgCO$_2$e，远高于基准碳价。
6）在现行时段结构下，单位碳价对充电时刻的作用存在门槛。取表\ref{tab:carbon-charging}中两种方案各10次运算求得的20组配送路线固定不变，
仅按电价与电网碳强度重排充电时刻，单位碳价自0以0.1元/kgCO$_2$e的增幅变化至3.0元/kgCO$_2$e：
碳价低于0.5元/kgCO$_2$e时无路线为减排多付电费，充电成本降幅保持40.07元不变，碳排量增加13.12 kgCO$_2$e；
0.5元/kgCO$_2$e起有3组路线为减排多付电费，0.8元/kgCO$_2$e起增至14组，碳排量增幅缩小到0.78 kgCO$_2$e；
只有当单位碳价超过2.0元/kgCO$_2$e后，考虑时变碳强度的充电安排才转为净减排（3.0元/kgCO$_2$e时减少4.60 kgCO$_2$e），
而只按碳强度择时可减排4.89 kgCO$_2$e。
上述结果表明，时变碳强度要成为值得企业与政府重视的信号，需要两个条件同时成立：分时电价谷段对齐到低碳时段，
单位碳价抬升到足以使最优车队电动化的水平。下一小节在此基础上比较不同碳减排政策的作用。

\subsubsection{不同碳减排政策对比分析}

文献中的碳减排政策工具按其对配送决策的作用可分为三类：
一是改变燃油车与电动车相对经济性、从而改变车队构成与车型指派的工具，如碳税或单位碳价\cite{ref:23,ref:qiu2024}与电动车购置补贴\cite{ref:wilson2026}；
二是改变时段内相对电价、只改变充电时刻的工具，如分时电价时段方案\cite{ref:93}、充电补贴\cite{ref:wu2022}与企业作业时间安排；
三是只平移成本、不进入决策的工具，如碳配额与交易，其配额与决策变量没有直接关系\cite{ref:lijin2014,ref:23}。
本文从三类中各取代表，设置表\ref{tab:policy-scenarios}所示的政策情形，其中电动车固定成本溢价下降对应电池价格与保险费率下降等市场因素，用于与购置补贴对照；
碳限额硬约束、碳抵消、绿电充电站与限行区等工具或需改变模型结构，或作用于本文未建模的空间维度，不在本文范围内。

\begin{table}[H]
  \centering
  \caption{碳减排政策情形设置}
  \label{tab:policy-scenarios}
  \small
  \begin{tabular*}{\textwidth}{@{\extracolsep{\fill}}llllc@{}}
    \toprule
    政策情形 & 作用途径 & 实施主体 & 取值 & 依据\\
    \midrule
    单位碳价 & 改变车队构成 & 政府 & 0.075（现行）、0.60、1.5、2.0元/kgCO$_2$e & \cite{ref:23,ref:qiu2024}\\
    电动车购置补贴 & 改变车队构成 & 政府 & 全额补贴购置价差，按使用年限折算为24元/日 & \cite{ref:wilson2026}\\
    电动车固定成本溢价下降 & 改变车队构成 & 市场 & 日固定成本溢价由100元降至50元、0元 & \cite{ref:wilson2026}\\
    碳配额与交易 & 平移成本 & 政府 & 配额0--500 kg，增幅100 kg & \cite{ref:lijin2014}\\
    分时电价谷段挪至午间 & 改变充电时刻 & 电网 & 谷段01:00--06:00与12:00--15:00 & \cite{ref:hebei-tou}\\
    午间充电按谷价补贴 & 改变充电时刻 & 政府 & 12:00--15:00充电按谷价计费 & \cite{ref:wu2022}\\
    午休时段后移 & 改变充电时刻 & 企业 & 午休由11:00--13:00改为11:00--14:00 & ---\\
    \bottomrule
  \end{tabular*}
\end{table}

为探究单位碳价对配送方案的影响，在北京现行分时电价下令单位碳价自0起以0.1元/kgCO$_2$e的增幅变化至2.0元/kgCO$_2$e，
另取现行碳价0.075元/kgCO$_2$e，各档独立运行3次；由于各档求得的配送方案在其他碳价下同样可行，
把全部66个方案作为方案池，在每一碳价下重新核算并取总成本最低者作为该碳价的最优配送方案，
对应方案的派遣电动车占比与总成本如图\ref{fig:carbon-price-sweep}所示。

%%FIGURE5%%

由图\ref{fig:carbon-price-sweep}可知：
1）派遣电动车占比随单位碳价呈阶梯式上升，总成本随之分段线性上升并在换挡处斜率变缓，阶梯只出现在两处：
单位碳价升至1.24元/kgCO$_2$e时最优车队由2辆燃油车3辆电动车改为1辆燃油车5辆电动车，电动车占比由60\%升至83.3\%；
升至1.52元/kgCO$_2$e时改为纯电动，
对应碳排量由191.55 kgCO$_2$e降至101.72 kgCO$_2$e再降至73.43 kgCO$_2$e；第一次翻转以增加111.53元非碳成本换取减排89.83 kgCO$_2$e，
第二次以增加43.07元换取减排28.29 kgCO$_2$e，减排的边际成本由1.24元/kgCO$_2$e升至1.52元/kgCO$_2$e。
2）现行碳价0.075元/kgCO$_2$e与基准碳价0.20元/kgCO$_2$e均落在第一段内，车队构型不变，碳价只改变碳成本的核算，
与陈婉茹等\cite{ref:23}报告的电动车占比在碳价0.5元/kg处才首次上升、Qiu等\cite{ref:qiu2024}报告的碳税低于0.1元/kg时排放几乎不动相一致。

为探究其余政策情形对配送方案的影响，逐一改变一种政策并重新优化车型、路线与充电时刻；
碳价、补贴、溢价、配额与作业时间情形各独立运行3次，基准与谷段挪至午间的情形运行10次，各项结果取其均值，燃油/电动列为总成本最低方案的车队构型。
表\ref{tab:fleet-levels}中给出各情形下配送方案的车队构型、启动成本、行驶成本、充电成本、油耗成本、碳成本、总成本、碳排量及其相对基准的变化量，
上半部为改变了车队构成的情形，下半部为车队不变、只改变充电时刻的情形。
各情形的碳成本随单位碳价与配额变化，比较总成本时应扣除该列。

%%TABLE13%%

各情形均完成50个客户和13264 kg需求。由表\ref{tab:fleet-levels}可知：
1）改变油电相对经济性的政策改变了车队构成。碳价在0.075--0.60元/kgCO$_2$e区间内最优配送方案仍为2辆燃油车3辆电动车；
升至1.5元/kgCO$_2$e时改为纯电动，碳排量减少了123.92 kgCO$_2$e（63.6\%），总成本增加200.32元（7.60\%）；
2.0元/kgCO$_2$e时碳排量减少了130.49 kgCO$_2$e（66.9\%），总成本增加238.51元（9.05\%）。
购置补贴按使用年限折算为24元/日后，最优车队改为2辆燃油车4辆电动车，总成本下降66.12元（2.51\%），碳排量减少51.83 kgCO$_2$e（26.6\%）；
电动车固定成本溢价降至50元/日时车队即转为纯电动，总成本下降199.41元，碳排量减少87.11 kgCO$_2$e。
这一门槛与文献一致：Qiu等\cite{ref:qiu2024}报告碳税低于0.1元/kg时排放几乎不动、0.1--1元/kg区间才明显下降，
陈婉茹等\cite{ref:23}报告电动车占比在碳价0.5、2.2、3.4、5.0元/kg处阶梯式上升。
2）碳配额与交易不改变配送方案。配额由0增至200 kg时，最优配送方案的路线、车队构型与趟数均不变，各项非碳成本逐位相同，
总成本仅按碳价与配额的乘积平移40.00元，碳成本转为负值即企业出售剩余配额获得收益；
在基准路线上把配额自0以100 kg的增幅变化至500 kg，全部路线的充电时刻逐位相同。
这是因为碳交易成本对排放量为线性，配额只是常数项，与李进等\cite{ref:lijin2014}和陈婉茹等\cite{ref:23}的结论一致。
3）分时电价时段方案与充电补贴只改变充电时刻。谷段挪至午间后最优配送方案的车队构型与基准同为2辆燃油车3辆电动车，
充电成本由172.29元降至129.53元，总成本下降39.20元（1.49\%），碳排量减少10.60 kgCO$_2$e；
午间充电按谷价补贴使充电成本降至140.76元，总成本下降54.31元（2.06\%），碳排量几乎相同，财政每日支出45.90元，属于转移支付。
这是因为在本算例的作业结构下，全天的充电只发生在开工前、午休回场与下午班趟间三处，
前两处的时刻由班次锁定，第三处由3个下午客户的时间窗锁定，其中2个客户最早回场已在谷段结束之后，
任何路线都无法把这部分充电移入谷段；按其他省份的实际谷段方案（10:00--15:00、11:00--16:00、12:00--14:00等）逐一核算，
路线改动可再多获得的电费收益均不足0.4元/日。可见分时电价时段方案的作用取决于谷段的位置，而不取决于路线的安排。
4）作业时间调整的作用方向不定。午休后移至11:00--14:00后充电成本下降34.18元，
但最优方案改用4辆燃油车2辆电动车，油耗成本上升51.72元，总成本与基准几乎相同，碳排量反而增加12.19 kgCO$_2$e，
即在车队油电经济性接近平价时，作业时间的变动可能把车型指派推向任一方向。

上述结果表明，改变车队构成的政策与改变充电时刻的政策作用于配送决策的不同维度。
为探究两类政策的组合作用，把分时电价谷段挪至午间与单位碳价升至1.00元/kgCO$_2$e两个因素组成四种情形，
每种情形仍按三种充电安排求解，即充方案与路线重优化方案的结果均基于10次运算，
对应方案的总成本与碳排量如表\ref{tab:grid2x2}所示。
其中，路线不变方案的车队构型与即充方案相同；本文将总成本减去碳成本在此称为运营成本。

%%TABLE12%%

由表\ref{tab:grid2x2}可知：
1）考虑时变碳强度在四种情形下均降低了成本，但只在两因素组合情形下同时降低了排放。
路线不变方案相对即充方案的总成本在四种情形下全部下降（降幅32.30--73.64元），
碳排量在基准情形与单独对齐时段的情形下分别增加11.28与10.21 kgCO$_2$e，单独抬高碳价时几乎相同（增加1.16 kgCO$_2$e），
只有谷段挪至午间且单位碳价为1.00元/kgCO$_2$e时减少4.38 kgCO$_2$e，10组路线全部同时下降。
2）单独对齐时段或单独抬高碳价，路线重优化都没有带来额外收益。只对齐时段时，路线重优化方案与路线不变方案的总成本相差10.24元，
碳排量比基准情形低9.86 kgCO$_2$e；只抬碳价时，两者总成本相差5.85元，路线重优化方案的碳排量比基准情形低41.80 kgCO$_2$e，
减排来自碳价推动的车队电动化（10次中5次为2辆燃油车4辆电动车），与充电时刻无关。
3）两者组合后，碳强度推动了车队构成改变。谷段挪至午间且单位碳价为1.00元/kgCO$_2$e时，
路线重优化方案10次运算中有9次派出4辆及以上电动车，即充方案为6次；三种安排的排序首次成立：
路线重优化方案的总成本比路线不变方案低44.67元（1.61\%）、碳排量低36.02 kgCO$_2$e（23.33\%），
比即充方案低118.31元（4.15\%）、40.39 kgCO$_2$e（25.44\%）。
4）与基准情形相比，两因素组合情形下路线重优化方案的碳排量减少了76.58 kgCO$_2$e（39.28\%），运营成本增加20.61元（0.79\%），
含碳支出的总成本增加100.00元，即企业多付的主要是碳价上升带来的碳支出本身。
运营成本几乎不变的原因在于车队电动化的两笔账相抵：油耗成本减少189.73元、行驶成本减少41.66元、充电成本减少16.99元，合计减少248.38元，
而启动成本因多派电动车增加269.00元，电动车每日100元的固定成本溢价与其能耗节省在本算例参数下接近平价；
碳价的作用是使企业跨过这一溢价门槛，减排主体为燃油车直接排放，由123.22 kgCO$_2$e降至56.20 kgCO$_2$e。
5）两种方案的趟数分布几乎相同，每次运算均为15或16趟，配送方案的改变体现在车型指派而非趟次结构上。

综上所述，对配送企业而言，在现行碳价与分时电价下按电价与碳强度安排充电只能降低电费，首趟前补电会被谷价引向碳强度最高的夜间时段，
真正的减排来自午休回场补电进入午后低碳时段；车队构成应按油电相对经济性决定，碳价与购置补贴只在跨过溢价门槛时才改变车队。
对政府而言，单位碳价只能初步推进配送车队的电动化，进一步推广还需把分时电价谷段对齐到光伏出力较高的午间，并辅以车辆购置支持；
碳配额的高低不改变企业的配送决策，起作用的是碳价。
上述管理启示也与陈婉茹等\cite{ref:23}提出的``限额碳交易政策只能初步推进新能源车辆在物流配送中的使用''的研究结论相吻合。

